\scanpage{4}
% archive foliation: 3
\deleted{\uline{Prop}}\inserted{\uline{Th\unsure{m} 1}}
\[
(C(j)\setminus C_0)\cap G(j)=\{x\in C(j)\inserted{\setminus C_0}\mid\deleted{\unsure{x\in G_0(j)}\ \unsure{\text{ou}}}\ \exists y\in R(x),\ \unsure{\text{avec}}\ y\in G(j)\}
\]
\inserted{$R(x)\cap G(j)\ne\varnothing$}
\[
(\overline C(j)\setminus C_0)\cap G(j)=\{x\in\overline C(j)\mid\forall y\in R(x)\ \text{on a }y\in G(j)\ \deleted{\text{et }\unsure{\text{si}}\ R(x)=\varnothing\ (\text{i.e. }x\in C_0)\ \unsure{\text{alors}}\ x\in G_0(j)}\}
\]
\inserted{$R(x)\subset G(j)$}
\margin{\unsure{S'il n'y a pas} de partie infinie}

\uline{Dém} \deleted{a) Soit $x\in C(j)$}\inserted{$1^\circ)$} Soit $x\in C(j)$, il faut prouver que
\[
x\in G(j)\Longleftrightarrow\deleted{\unsure{x\in G_0(j)\ \text{ou}}}\ \exists y\in R(x),\ y\in G(j).
\]
\deleted{\unsure{x est terminal}}

\deleted{Si $x\in C_0$, c'est clair (cf NB plus haut)}
Soit $x\in C\setminus C_0$ \inserted{i.e. $R(x)\ne0$. Prouvons $\Longrightarrow$ i.e. supposons $x\in G(j)$, prouvons $\exists y$ \unsure{ch}\ldots}, soit $\Sigma$ stratégie \deleted{pour $j$} $x$-gagnante. Considérons $\Sigma(x)$; \unsure{donc} \deleted{nécessairement} $\Sigma(x)\ne\varnothing$, donc $\exists y\in\Sigma(x)$. Je dis que $y\in G(j)$, i.e.

\uline{Lemme} Si $x\in C(j)\cap G(j)$, $\Sigma$ une $j$-stratégie $x$-gagnante, alors $\forall y\in\Sigma(x)$, \inserted{on a} $y\in G(j)$, \inserted{i.e. $\Sigma(x)\subset G(j)$.} Plus précisément $\Sigma$ est aussi $y$-gagnante.

\deleted{En effet, \unsure{parties pour} $\Sigma$}
Trivial sur définitions.

Inversement, supposons $\exists y\in R(x)$ avec $y\in G(j)$, prouvons que $x\in G(j)$. Soit $\Sigma$ une $j$-stratégie $y$-gagnante. On va \inserted{\unsure{en garder|en gardant}} définir une $j$-stratégie \inserted{$\Sigma'$} $x$-gagnante ainsi : Pour $z\in C(j)$, on pose
\[
\Sigma'(z)=\begin{cases}\Sigma(z)\text{ si }z\ne x\\ \{y\}\deleted{\unsure{\text{ou}}}\text{ si }z=x\end{cases}
\]
C'est une $j$-stratégie, car \deleted{si $R(z)\ne\varnothing$ \unsure{-} $R(x)\ne\varnothing$, alors} $\Sigma(z)\ne0$ \deleted{\unsure{pour}} \deleted{(et dans le cas $z\ne x$ pas de pb, mais alors $\Sigma'(z)=\Sigma(z)\ne\varnothing$ par hyp. que $\Sigma$ est stratégie.}

Prouvons que $\Sigma'$ est $x$-gagnante : Soit une partie $(x_0=x,x_1,\ldots,x_n\ldots)$ partant de $x$ et compatible avec $\Sigma'$. On a donc $x_1=y$; \deleted{et}\inserted{prouvons que} la partie $(x_1,x_2,\ldots,x_n\ldots)$ \inserted{$x_1=y$} est compatible avec $\Sigma$, i.e. que
\[
\deleted{\unsure{\forall}}\quad i\geq1,\quad x_i\in C(j)\Longrightarrow x_{i+1}\in\Sigma(x_i).
\]
